\section{Existence of continuous equivalent layers on a plane}
\label{sec:planar_eqlayer_theorem}

\subsection{Modified Dirichlet's and Neumann's problems on a plane}

To prove the existence of continuous equivalent layers $p_{\varepsilon}(\mathbf{r}'')$
according to Definition \ref{def:continuous-eqlayer},
we need first modify the Dirichlet’s and Neumann’s problems 
(equations \ref{eq:U-Dirichlet-Neumann} and \ref{eq:U-Dirichlet-Neumann-Fourier}) for the case in which
the plane $z''_{0}$ cross the sources.
Hence, let us first consider an approximating potential $U_{\varepsilon}(\mathbf{r}'')$ 
defined on a horizontal plane $z''_{0}$ that may be located above or cross 
the sources (Figure \ref{fig:geometry}), with Fourier transform given by
\begin{equation}
	\mathcal{F} \left[ U_{\varepsilon}(\mathbf{r}'') \right] = 
	\begin{cases}
		\mathcal{F} \left[ U(\mathbf{r}^{\dagger}) \right] 
		e^{\, k \, (z''_{0} - z^{\dagger}_{0})} & , 
		\quad 0 < k \le k_{\varepsilon} \\
		0 & , \quad k > k_{\varepsilon} 
	\end{cases} \: ,
	\label{eq:potential-U-epsilon-Fourier-layer}
\end{equation}
where $k_{\varepsilon}$ is a fixed and possibly very large truncating wavenumber and $z^{\dagger}_{0} < z^{t}_{0}$. 
Consider also that the inverse Fourier transform of $\mathcal{F} \left[ U_{\varepsilon}(\mathbf{r}'') \right]$ 
(equation \ref{eq:potential-U-epsilon-Fourier-layer}) satisfies
\begin{equation}
	U_{\varepsilon}(\mathbf{r}) = 
	\int\limits_{-\infty}^{\infty} \int\limits_{-\infty}^{\infty}
	D(\ell, \hat{\mathbf{z}}) U_{\varepsilon}(\mathbf{r}'') \:\:
	D(1-\ell, \hat{\mathbf{z}}) \Psi(\mathbf{r}, \mathbf{r}'') 
	\:\: d\mathcal{S}'' \: , 
	\quad \forall z < z^{t}_{0} \: , \quad \forall z < z''_{0} \: ,
	\quad \ell \in \{ 0, 1 \} \: ,
	\label{eq:potential-U-epsilon-space}
\end{equation}
where $D$ is the space-domain differential operator 
defined by equation \ref{eq:diff-operators-generic} in terms of the canonical base vector 
$\hat{\mathbf{z}}$ (equation \ref{eq:canonical-base-vector}).

Note that the function $U_{\varepsilon}(\mathbf{r}'')$ obtained by computing
the inverse Fourier transform of 
$\mathcal{F} \left[ U_{\varepsilon}(\mathbf{r}'') \right]$ (equation \ref{eq:potential-U-epsilon-Fourier-layer})
approximates the potential field obtained by 
computing the downward continuation of $U$ (equation \ref{eq:potential-U-space}) 
to points that may be located inside the sources. 
It is well-known, however, that downward continuation can only be computed at 
source-free points, where the Laplace's equation applies. This means that the
continuation plane is limited by the top $z^{t}_{0}$ of the sources, otherwise the continued field 
presents an oscillatory character \citep[e.g.,][]{roy1967, fedi_florio2002} that may
prevent $\mathcal{F} \left[ U_{\varepsilon}(\mathbf{r}'') \right]$ to be square integrable.
The truncating wavenumber $k_{\varepsilon}$ (equation \ref{eq:potential-U-epsilon-Fourier-layer}) 
overcomes this problem, as suggested by \citet{bhattacharyya1967},
and guarantees that the Fourier-convolution theorem 
can be applied to the auxiliary potential $U_{\varepsilon}$
(equation \ref{eq:potential-U-epsilon-space}) on a horizontal plane $z^{\dagger}_{0} < z^{t}_{0}$ 
to obtain
\begin{equation}
	\mathcal{F} [ U_{\varepsilon}(\mathbf{r}^{\dagger}) ]
	= k \, \mathcal{F} [ U_{\varepsilon}(\mathbf{r}'') ] \:
	\mathcal{F} [ \Psi(\mathbf{r}^{\dagger}, \mathbf{r}'') ] \: .
	\label{eq:potential-U-epsilon-Fourier-dagger}
\end{equation}
We get by substituting \ref{eq:Psi-Fourier} and \ref{eq:potential-U-epsilon-Fourier-layer} into 
\ref{eq:potential-U-epsilon-Fourier-dagger} that
\begin{equation}
	\mathcal{F} [ U_{\varepsilon}(\mathbf{r}^{\dagger}) ] = 
	\begin{cases}
		\mathcal{F} [ U(\mathbf{r}^{\dagger}) ] &
		, \quad 0 < k \le k_{\varepsilon} \\
		0 & , \quad k > k_{\varepsilon} 
	\end{cases} \: ,
	\label{eq:potential-U-epsilon-Fourier-dagger2}
\end{equation}
which shows that the Fourier transforms of 
$U(\mathbf{r}^{\dagger})$ (equation \ref{eq:potential-U-Fourier}) and 
$U_{\varepsilon}(\mathbf{r}^{\dagger})$ (equation \ref{eq:potential-U-epsilon-space}) 
on the horizontal plane $z^{\dagger}_{0}$ located above the 3D sources coincide up 
to the truncating wavenumber $k_{\varepsilon}$.

We have just specified the Dirichlet's and Neumann's problems 
(equations \ref{eq:U-Dirichlet-Neumann} and \ref{eq:U-Dirichlet-Neumann-Fourier}) in terms of a truncating wavenumber 
$k_{\varepsilon}$ with the purpose of expanding their applicability to the case in which plane $z''_{0}$ cross the sources. 
In this case, we obtain an approximating potential $U_{\varepsilon}(\mathbf{r})$ that coincides with the true 
potential $U(\mathbf{r})$ up to the truncating wavenumber $k_{\varepsilon}$ 
(equation \ref{eq:potential-U-epsilon-Fourier-dagger2}). 
At last, it can be easily verified that, if we let the truncating wavenumber $k_{\varepsilon}$ 
goes to infinite in equation \ref{eq:potential-U-epsilon-Fourier-layer} for the case in which $z''_{0} < z^{t}_{0}$, we get 
\begin{equation}
	\lim_{k_{\varepsilon}\to\infty} \mathcal{F} \left[ U_{\varepsilon}(\mathbf{r}'') \right] = 
	\mathcal{F} \left[ U(\mathbf{r}'') \right] \: , \quad z''_{0} < z^{t}_{0} \: ,
	\label{eq:lim-kepsilon-potential-U-epsilon-Fourier-layer}
\end{equation}
\begin{equation}
	\lim_{k_{\varepsilon}\to\infty} \mathcal{F} \left[ U_{\varepsilon}(\mathbf{r}^{\dagger}) \right] = 
	\mathcal{F} \left[ U(\mathbf{r}^{\dagger}) \right] \: , \quad z^{\dagger}_{0} < z''_{0} < z^{t}_{0} \: ,
	\label{eq:lim-kepsilon-potential-U-epsilon-Fourier-dagger}
\end{equation}
pointwise and, consequently, that
\begin{equation}
	\lim_{k_{\varepsilon}\to\infty} U_{\varepsilon}(\mathbf{r}) = 
	U(\mathbf{r}) \: , \quad \forall z < z''_{0} < z^{t}_{0} \: ,
	\label{eq:lim-kepsilon-potential-U-epsilon-space}
\end{equation}
which is in perfect agreement with the solutions of the Dirichlet's and Neumann's
problems (equations \ref{eq:U-Dirichlet-Neumann} and \ref{eq:U-Dirichlet-Neumann-Fourier}) on a plane $z''_{0}$ 
that is located above the sources.

\subsection{The planar equivalent-layer theorem}

\begin{theorem}[\textbf{Planar equivalent-layer theorem}]
	\label{theorem:planar-eqlayer-theorem}
	Let a planar surface $\mathcal{S}''$ located at $z''_{0}$ and a kernel function $a(\mathbf{r}, \mathbf{r}'')$ 
	be set according to Definition \ref{def:continuous-eqlayer}, for a given observed potential field $d^{o}$. 
	Then, for every $\varepsilon > 0$, there exists a continuous equivalent layer $p_{\varepsilon}$. Furthermore, if 
	$z''_{0} < z^{t}_{0}$, then the continuous equivalent layer reproduces $d^{o}$ exactly.
\end{theorem}
\begin{proof}
	We start by defining $p_{\varepsilon}$ (Definition \ref{def:continuous-eqlayer}) 
	in the Fourier domain via
	\begin{equation}
		\mathcal{F} [ p_{\varepsilon}(\mathbf{r}'') ] =
		\tfrac{k \, \alpha_{d}}{\alpha_{a}} \,
		\mathcal{F} [ U_{\varepsilon}(\mathbf{r}'') ] \: ,
		\label{eq:p-Fourier}
	\end{equation}
	where $U_{\varepsilon}(\mathbf{r}'')$ is defined as in \ref{eq:potential-U-epsilon-Fourier-layer}, 
	for a given truncating wavenumber $k_{\varepsilon}$, and $\alpha_{a}$ is differential operator 
	defined in the Fourier domain in terms of $\alpha$ (equation \ref{eq:diff-operators-generic}).
	For any finite choice of $0 < k_{\varepsilon} < \infty$, both $\mathcal{F} [ p_{\varepsilon}(\mathbf{r}'') ]$ 
	and $p_{\varepsilon}(\mathbf{r}'')$ are square-integrable and, consequently, it is possible to
	apply the Fourier-convolution theorem to \ref{eq:predicted-data-continuous-layer} and obtain
	\begin{equation}
		\mathcal{F} [ d_{\varepsilon}(\mathbf{r}^{\dagger}) ] =
		\mathcal{F} [ p_{\varepsilon}(\mathbf{r}'') ] \, 
		\mathcal{F} [ a(\mathbf{r}^{\dagger}, \mathbf{r}'') ] \: ,
		\quad z^{\dagger}_{0} < z''_{0} \: ,
		\label{eq:predicted-data-continuous-layer-Fourier}
	\end{equation}
	where 
	\begin{equation}
		\mathcal{F} [ a(\mathbf{r}^{\dagger}, \mathbf{r}'') ] = \alpha_{a} \, \mathcal{F} [ \Psi(\mathbf{r}^{\dagger}, \mathbf{r}'') ] \: ,
		\label{eq:kernel-a-Fourier}
	\end{equation}
	with $\mathbf{r}^{\dagger} = (x^{\dagger}, y^{\dagger}, z^{\dagger}_{0})^{\top} \notin \overline{\mathcal{V}}$ 
	defining a point on the horizontal surface $\mathcal{S}^{\dagger}$ above the sources, at $z^{\dagger}_{0} < z''_{0}$.
	By considering that the kernel function is given by $a(\mathbf{r}, \mathbf{r}'') = D_{a} \Psi(\mathbf{r}, \mathbf{r}'')$
	and using \ref{eq:observed-data-Fourier} and \ref{eq:Psi-Fourier}, we find that $d_{\varepsilon}$ satisfies
	\begin{equation}
		\mathcal{F} [ d_{\varepsilon}(\mathbf{r}^{\dagger}) ] = 
		\begin{cases}
			\mathcal{F} [ d^{o}(\mathbf{r}^{\dagger}) ] &
			, \quad 0 < k \le k_{\varepsilon} \\
			0 & , \quad k > k_{\varepsilon} 
		\end{cases} \: .
		\label{eq:predicted-data-continuous-layer-Fourier2}
	\end{equation}
	From Plancherel's theorem \citep[][ p. 119]{bracewell1999} and use of polar coordinates 
	$k_{x} = k \, \cos\theta$, $k_{y} = k \, \sin\theta$ and $k = \sqrt{k_{x}^{2} + k_{y}^{2}}$,
	we get
	\begin{equation}
		\begin{split}
		\int\limits_{-\infty}^{\infty} \int\limits_{-\infty}^{\infty}
		\vert 
		d^{o}(\mathbf{r}^{\dagger}) - d(\mathbf{r}^{\dagger})
		\vert^{2} 
		\: d\mathcal{S}^{\dagger} = 
		\int\limits_{-\infty}^{\infty} \int\limits_{-\infty}^{\infty}
		\vert \mathcal{F} [ d^{o}(\mathbf{r}^{\dagger}) ] - \mathcal{F} [
		d(\mathbf{r}^{\dagger}) ] \vert^{2} 
		\: d k_{x} \, d k_{y} \\
		= \int\limits_{0}^{2\pi} \int\limits_{0}^{\infty}
		\vert \mathcal{F} [ d^{o}(\mathbf{r}^{\dagger}) ] - \mathcal{F} [
		d(\mathbf{r}^{\dagger}) ] \vert^{2} \: k \, dkd\theta =
		\int\limits_{0}^{2\pi} \int\limits_{k_{\varepsilon}}^{\infty}
		\vert \mathcal{F} [ d^{o}(\mathbf{r}^{\dagger}) ] \vert^{2} 
		\: k \: d k \, d \theta \: .
		\end{split}
		\label{eq:Plancherel}
	\end{equation}
	Since $\mathcal{F} [ d^{o}(\mathbf{r}^{\dagger}) ] \, k$ is square integrable on $\mathcal{S}^{\dagger}$, we may choose 
	a finite but sufficiently large $k_{\varepsilon} < \infty$ such that 
	$\int_{-\infty}^{\infty} \int_{-\infty}^{\infty} \vert d^{o}(\mathbf{r}^{\dagger}) - d(\mathbf{r}^{\dagger}) 
	\vert^{2} \: d\mathcal{S}^{\dagger} = \int_{0}^{2\pi} \int_{0}^{\infty} 
	\vert \mathcal{F} [ d^{o}(\mathbf{r}^{\dagger}) ] - \mathcal{F} [d(\mathbf{r}^{\dagger}) ] \vert^{2} \: k \, dkd\theta 
	\le \varepsilon$. This leads to the conclusion that $p_{\varepsilon}$ is a continuous equivalent layer. 
	If the layer plane is above the sources (i.e., if $z''_{0} < z^{t}_{0}$), then both 
	$\mathcal{F} [ p_{\varepsilon}(\mathbf{r}'') ]$ and $p_{\varepsilon}(\mathbf{r}'')$ are square-integrable
	even for the choice $k_{\varepsilon} = \infty$. Then, in that case, 
	$\int_{-\infty}^{\infty} \int_{-\infty}^{\infty} \vert d^{o}(\mathbf{r}^{\dagger}) - d(\mathbf{r}^{\dagger}) 
	\vert^{2} \: d\mathcal{S}^{\dagger} = \int_{0}^{2\pi} \int_{0}^{\infty} 
	\vert \mathcal{F} [ d^{o}(\mathbf{r}^{\dagger}) ] - \mathcal{F} [d(\mathbf{r}^{\dagger}) ] \vert^{2} \: k \, dkd\theta 
	= 0$ and $p_{\varepsilon}$ yields a predicted field $d_{\varepsilon}$ that exactly reproduces the observed field $d^{o}$ 
	on $\mathcal{S}^{\dagger}$.
\end{proof}
	
\subsubsection{Remark 1}

The above theorem guarantees the existence of an approximate equivalent layer, even if
the layer crosses the sources, i.e., if $z''_{0} > z^{t}_{0}$.
But it does not provide any guidance on the particular choice of the truncating wave number $k_{\varepsilon}$ 
for the construction of $p_{\varepsilon}$.
However, it is possible to combine the upper bounds defined in Appendix A 
with the right-hand side of equation \ref{eq:Plancherel} to show that
\begin{equation}
	\int\limits_{0}^{2\pi} \int\limits_{0}^{\infty} 
	\vert \mathcal{F} [ d^{o}(\mathbf{r}^{\dagger}) ] - \mathcal{F} [d(\mathbf{r}^{\dagger}) ] \vert^{2} \: k \, dkd\theta 
	\le \varepsilon(q, k_{\varepsilon}, z^{\dagger}_{0}) \: ,
	\label{eq:Plancherel-epsilon}
\end{equation}
where
\begin{equation}
	\varepsilon(q, k_{\varepsilon}, z^{\dagger}_{0}) = \tfrac{2\pi \, \Omega^{2} \, 3^{q}}{(2 \, \Delta z)^{2q}} \: 
	\Gamma(2 \, q \: , \: 2 \, \Delta z \, k_{\varepsilon}) \: ,
	\quad \Delta z = z^{t}_{0} - z^{\dagger}_{0} > 0 \: ,
	\label{eq:epsilon}
\end{equation}
is a real function depending on $\Omega$ (Apendix A, equation \ref{eq:Omega}) and the Incomplete Gamma function 
\citep[][ p. 899]{gradshteyn_ryzhik2007}
\begin{equation}
	\Gamma(2 \, q \: , \: 2 \, \Delta z \, k_{\varepsilon}) = 
	\int\limits_{k_{\varepsilon}}^{\infty}
	e^{- 2 \, \Delta z \, k } \, k^{2q-1} \, (2 \, \Delta z)^{2q} 
	\: d k \: ,
	\label{eq:incomplete-Gamma-function}
\end{equation}
with $q$ being the order of the differential operator $\alpha_{d}$ (equation \ref{eq:diff-operators-generic}). 
According to equation \ref{eq:epsilon}, if we have some a priori bounds 
on the physical-property distribution and position of the true sources, we get a connection between the choice of 
$k_{\varepsilon}$ and the desired accuracy level $\varepsilon$. 
This can potentially support an estimation of the truncating wave number $k_{\varepsilon}$.

\subsubsection{Remark 2}

The truncating wavenumber $k_{\varepsilon}$ plays a central role in the planar equivalent-layer theorem 
(Theorem \ref{theorem:planar-eqlayer-theorem}).
We would like to highlight here that the interpreter typically does not explicitly set $k_{\varepsilon}$ in
practical situations. In fact, obtaining a discrete equivalent layer (Definition \ref{def:discrete-eqlayer}) 
does not strictly imply a truncation of the wavenumber at some maximal $k_{\varepsilon}$, 
but it leads to an exponentially fast decay at sufficiently large wavenumbers. 
From a theoretical point of view, it might be interesting to study the influence of
such a discretization and its relation to $k_{\varepsilon}$ in more detail. 
For practical purposes it seems that this discretization is the predominant factor rendering 
the equivalent-layer technique applicably even when the layer plane crosses the sources.

\subsubsection{Remark 3}

The equivalent layer from Theorem \ref{theorem:planar-eqlayer-theorem} leads to a predicted 
potential field that does not satisfy
the error bound \ref{eq:epsilon-bound-data} only on the surface $\mathcal{S}^{\dagger}$, but also 
on any $\mathcal{S}^{\ddagger}$ with $z^{\ddagger}_{0} < z^{\dagger}_{0}$. 
Namely, it follows from equation \ref{eq:potential-U-Fourier} that
\begin{equation}
	\int\limits_{-\infty}^{\infty} \int\limits_{-\infty}^{\infty}
	\vert d^{o}(\mathbf{r}^{\ddagger}) - d(\mathbf{r}^{\ddagger})
	\vert^{2} \: d\mathcal{S}^{\ddagger} \le
	\int\limits_{-\infty}^{\infty} \int\limits_{-\infty}^{\infty}
	\vert d^{o}(\mathbf{r}^{\dagger}) - d(\mathbf{r}^{\dagger})
	\vert^{2} \: d\mathcal{S}^{\dagger} \: ,
	\quad z^{\ddagger}_{0} < z^{\dagger}_{0} \: .
	\label{eq:error-bound-upper-planes}
\end{equation}
Thus, if $p_{\varepsilon}$ produces a $d_{\varepsilon}$ satisfying 
$\int_{-\infty}^{\infty} \int_{-\infty}^{\infty} \vert d^{o}(\mathbf{r}^{\dagger}) - d(\mathbf{r}^{\dagger}) 
\vert^{2} \: d\mathcal{S}^{\dagger} \le \varepsilon$, then it automatically holds that 
$\int_{-\infty}^{\infty} \int_{-\infty}^{\infty} \vert d^{o}(\mathbf{r}^{\ddagger}) - d(\mathbf{r}^{\ddagger}) 
\vert^{2} \: d\mathcal{S}^{\ddagger} \le \varepsilon$.

\section{Potential-field transformations via equivalent-layer technique}

Note that, given a particular equivalent layer $\tilde{p}_{\varepsilon}$ (Definition \ref{def:continuous-eqlayer}) 
defined in terms of specific $\alpha_{d}$ and $\alpha_{a}$ as a square-integrable function by equation \ref{eq:p-Fourier}, 
it is theoretically possible 
to compute any potential field transformation (equation \ref{eq:transformed-field-space} and \ref{eq:transformed-field-Fourier}).
This can be verified by first noting that applying the Fourier-convolution theorem to \ref{eq:transformed-field-continuous-layer} 
on a planar surface $\mathcal{S}^{\ddagger}$ at a constant vertical coordinate $z^{\ddagger}_{0} < z''_{0}$ gives
\begin{equation}
	\mathcal{F} [ f_{\varepsilon}(\mathbf{r}^{\ddagger}) ] =
	\mathcal{F} \left[ \tilde{p}_{\varepsilon}(\mathbf{r}'') \right]
	\mathcal{F} [ t(\mathbf{r}^{\ddagger}, \mathbf{r}'') ] \: , 
	\quad z^{\ddagger}_{0} < z''_{0} \: ,
	\label{eq:transformed-field-continuous-layer-Fourier}
\end{equation}
where
\begin{equation}
	\mathcal{F} [ t(\mathbf{r}^{\ddagger}, \mathbf{r}'') ] = \alpha_{t} \, \mathcal{F} [ \Psi(\mathbf{r}^{\ddagger}, \mathbf{r}'') ]
	\label{eq:kernel-t-Fourier}
\end{equation}
is defined in terms of a differential operator $\alpha_{t}$ acting on $\Psi$ (equation \ref{eq:Psi}) in the Fourier domain.
Then we set 
\begin{equation}
	\alpha_{t} = \left( \tfrac{\alpha_{a}}{\alpha_{d}} \right) \alpha_{f} \: ,
	\label{eq:alpha-t-generic}
\end{equation}
where $\alpha_{f}$ is the differential operator associated with the desired transformed field $f$ (equation \ref{eq:transformed-field-Fourier}).
In this case, we get from equations \ref{eq:Psi-Fourier}, \ref{eq:p-Fourier}  and \ref{eq:kernel-t-Fourier} 
that the transformed potential field $f_{\varepsilon}$ (equation \ref{eq:transformed-field-continuous-layer-Fourier}) 
is given by:
\begin{equation}
	\mathcal{F} [ f_{\varepsilon}(\mathbf{r}^{\ddagger}) ] =
	\alpha_{f} \, k \, \mathcal{F} \left[ U_{\varepsilon}(\mathbf{r}'') \right]
	\mathcal{F} [ \Psi(\mathbf{r}^{\ddagger}, \mathbf{r}'') ] \: , 
	\quad z^{\ddagger}_{0} < z''_{0} \: .
	\label{eq:transformed-field-continuous-layer-Fourier2}
\end{equation}
Finally, we use \ref{eq:Psi-Fourier}, \ref{eq:transformed-field-Fourier} and \ref{eq:potential-U-epsilon-Fourier-layer} to show that
\begin{equation}
	\mathcal{F} [ f_{\varepsilon}(\mathbf{r}^{\ddagger}) ] = 
	\begin{cases}
		\mathcal{F} [ f(\mathbf{r}^{\ddagger}) ] &
		, \quad 0 < k \le k_{\varepsilon} \\
		0 & , \quad k > k_{\varepsilon} 
	\end{cases} \: ,
	\label{eq:transformed-field-continuous-layer-Fourier3}
\end{equation}
which shows that the transformed field $f_{\varepsilon}$ yielded by the equivalent layer coincides with 
that produced by the true sources on the planar surface $\mathcal{S}^{\ddagger}$ up 
to the truncating wavenumber $k_{\varepsilon}$.

In the context of the equivalent-layer technique, a potential-field transformation is generally defined
in the space domain, by following two steps. The first consists in obtaining a particular layer $\tilde{p}_{\varepsilon}$ 
(Definition \ref{def:continuous-eqlayer}), which is subsequently used to compute a transformed field $f_{\varepsilon}$ via 
\ref{eq:transformed-field-continuous-layer}.
It can be verified, however, that this two-step approach is equivalent to the single-step approach based on the 
Fourier transform of the observed potential field (equation \ref{eq:field-transformation-Fourier}). 
We find by using \ref{eq:observed-data-Fourier} to manipulate \ref{eq:transformed-field-continuous-layer-Fourier2} that
\begin{equation}
	\mathcal{F} [ f_{\varepsilon}(\mathbf{r}^{\ddagger}) ] =
	\left( \tfrac{\alpha_{f}}{\alpha_{d}} \right) \, 
	\begin{cases}
		\mathcal{F} \left[ d^{o}(\mathbf{r}^{\dagger}) \right] 
		e^{\, k \, (z^{\ddagger}_{0} - z^{\dagger}_{0})} & , \quad 0 < k \le k_{\varepsilon} \\
		0 & , \quad k > k_{\varepsilon} 
	\end{cases} \: .
	\label{eq:transformed-field-continuous-layer-single-step}
\end{equation}
Note that, for the particular case in which $z^{\ddagger}_{0}$ coincides with $z^{\dagger}_{0}$ and 
the truncating wavenumber $k_{\varepsilon}$ goes to infinite, equation 
\ref{eq:transformed-field-continuous-layer-single-step} reduces to \ref{eq:field-transformation-Fourier}.
This result establishes a theoretical link between potential-field transformations made via
equivalent-layer technique and those commonly made via Fourier transform of the observed data.
This link explains, for example, why potential-field transformations made via equivalent-layer technique may suffer from the same 
numerical problems generally associated with those made via Fourier transform of the observed potential-field data, 
as have already been pointed out by \citet{li_etal2014} about the reduction to the pole.

\section{Equivalent layer classes}
\label{sec:eqlayer_classes}

\subsection{Dirichlet's equivalent layers}

\begin{definition}
	An equivalent layer is called Dirichlet's equivalent layer if the Fourier transform of its physical-property distribution has 
	the form
	\begin{equation}
		\mathcal{F} [ p(\mathbf{r}'') ] =
		\begin{cases}
			\alpha_{d} \, \mathcal{F} [ U(\mathbf{r}^{\dagger}) ] 
			e^{\, k \, (z''_{0} - z^{\dagger}_{0})} & , 
			\quad 0 < k \le k_{\varepsilon} \\
			0 & , \quad k > k_{\varepsilon} 
		\end{cases} \: .
		\label{eq:p-Fourier-Dirichlet}
	\end{equation}
	where $z''_{0} > z^{\dagger}_{0}$.
	\label{def:eql-Dirichlet}
\end{definition}

Note that the particular $\mathcal{F} [ p(\mathbf{r}'') ]$ associated with Dirichlet's layers (equation \ref{eq:p-Fourier-Dirichlet}) 
is defined from the general solution (equation \ref{eq:p-Fourier}) by setting $\alpha_{a} \equiv \alpha(1, \hat{\mathbf{z}})$ 
(equations \ref{eq:diff-operators-generic}, \ref{eq:canonical-base-vector} and \ref{eq:kernel-Fourier}).
It is important noting that this particular $\mathcal{F} [ p(\mathbf{r}'') ]$ depends on $\alpha_{d}$, which 
carries information about the observed potential field (equation \ref{eq:observed-data-Fourier}).
In this case, the planar equivalent-layer integral in the space domain (equation \ref{eq:eqlayer-integral}) assumes the form
\begin{equation}
	d^{o}(\mathbf{r}) \approx
	d(\mathbf{r})
	= \int\limits_{-\infty}^{\infty}\int\limits_{-\infty}^{\infty}
	p(\mathbf{r}'') \: \tfrac{\partial}{\partial z} \Psi(\mathbf{r}, \mathbf{r}'') \:\: d\mathcal{S}''
	\: , \quad z < z''_{0} \: .
	\label{eq:eqlayer-integral-Dirichlet-space}
\end{equation}
Note that the solution $p(\mathbf{r}'')$ obtained by computing the inverse Fourier transform of 
$\mathcal{F} [ p(\mathbf{r}'') ]$ (equation \ref{eq:p-Fourier-Dirichlet}) approximates the
downward continuation of the observed potential field to the layer plane.
From the physical point of view, Dirichlet's layers give rise to a predicted potential field $d(\mathbf{r})$ which is analogous
to the vertical component of the gravitational attraction due to a layer of monopoles
or the magnetic scalar potential due to a layer of vertically-magnetized dipoles.
Moreover, it can also be verified that in this case the planar equivalent-layer integral 
(equation \ref{eq:eqlayer-integral}) approximates the well-known \textit{upward continuation integral} 
\citep[e.g.,][]{skeels1947, peters1949, henderson-zietz1949, dean1958, roy1962, bhattacharyya1967, henderson1970},
which is traditionally derived from the solution of the Dirichlet’s problem on the
layer plane (equation \ref{eq:U-Dirichlet-Neumann} with $\ell = 0$). 
Dirichlet's equivalent layers are the basis of the pioneer method proposed by \cite{dampney1969} for processing gravity data.

The particular physical-property distribution $\tilde{p}(\mathbf{r}'')$ satisfying \ref{eq:eqlayer-integral-Dirichlet-space}
for a given observed potential field $d^{o}(\mathbf{r})$ leads to potential-field transformations (equations 
\ref{eq:transformed-field-continuous-layer} and \ref{eq:transformed-field-integral-Fourier}) defined by 
$\alpha_{t} = \nicefrac{k \, \alpha_{f}}{\alpha_{d}}$ (equation \ref{eq:alpha-t-generic}).
Dirichlet's equivalent layers are commonly used to perform upward/downward continuation of the observed potential 
field. In this case, $\alpha_{f} = \alpha_{d}$ and $\alpha_{t} = k$. 
Other potential-field transformations can also be accomplished from such equivalent layers, but they must be 
defined in terms of $\alpha_{f} \ne \alpha_{d}$ and $\alpha_{t} \ne k$.
In this case, however, defining the potential-field transformation in the space domain (equation \ref{eq:transformed-field-continuous-layer})
from its counterpart in the Fourier domain (equation \ref{eq:transformed-field-integral-Fourier}) is generally 
non-trivial.


\subsection{Equivalent layers having minimum phase changes}

\begin{definition}
	An equivalent layer has minimum phase changes (with respect to the spatial distribution of the true sources) if its
	physical-property distribution has a phase spectrum given by 
	\begin{equation}
		\Theta [ p(\mathbf{r}'') ] = \begin{cases}
			\Theta [ U(\mathbf{r}^{\dagger}) ] & , \quad 0 < k \le k_{\varepsilon} \\
			0 & , \quad k > k_{\varepsilon}
		\end{cases} \: ,
		\label{eq:p-phase-spectrum-U}
	\end{equation}
	where $z''_{0} > z^{\dagger}_{0}$ and $\Theta [ U(\mathbf{r}^{\dagger}) ]$ (Appendix B, equation \ref{eq:potential-U-phase-spectrum})
	is the phase spectrum of the potential $U(\mathbf{r})$ (equation \ref{eq:potential-U-space}).
	\label{def:eql-min-phase}
\end{definition}
\begin{corollary}
	\label{corollary:minimum-phase-change-layers}
	An equivalent layer having minimum phase changes (Definition \ref{def:eql-min-phase})
	depends only on the spatial distribution of the true sources and not on the potential field
	used to define the layer.
\end{corollary}
\begin{proof}
	It follows from equations \ref{eq:p-phase-spectrum-U} and \ref{eq:potential-U-phase-spectrum} that 
	$\mathcal{F} [ p(\mathbf{r}'') ]$ must be the product of a real function and
	$\mathcal{F} [ U(\mathbf{r}^{\dagger}) ]$, for $0 < k \le k_{\varepsilon}$, or $0$, for $k > k_{\varepsilon}$.
	From equations \ref{eq:potential-U-epsilon-Fourier-layer} and \ref{eq:p-Fourier}, we conclude that this condition is satisfied only for the case in which 
	$\alpha_{a} = \alpha_{r} \, \alpha_{d}$, where $\alpha_{r}$ is a real function satisfying 
	$\left| \alpha_{r} \right| \ne 0$ for $k > 0$.
	In this case, it can be verified 
	that $\mathcal{F} [ a(\mathbf{r}^{\dagger}, \mathbf{r}'') ]$ (equation \ref{eq:kernel-Fourier}) 
	assumes the form
	\begin{equation}
		\mathcal{F} [ a(\mathbf{r}^{\dagger}, \mathbf{r}'') ] =
		\alpha_{r} \, \alpha_{d} \, 
		\mathcal{F} [ \Psi(\mathbf{r}^{\dagger}, \mathbf{r}'') ] \: ,
		\quad \left| \alpha_{r} \right| \ne 0 \: ,
		\label{eq:kernel-Fourier-min-phase}
	\end{equation}
	where $\mathcal{F} [ \Psi(\mathbf{r}^{\dagger}, \mathbf{r}'') ]$ is given by equation \ref{eq:Psi-Fourier} 
	and $\mathcal{F} [ p(\mathbf{r}'') ]$ (equation \ref{eq:p-Fourier}) assumes the form
	\begin{equation}
		\mathcal{F} [ p(\mathbf{r}'') ] =
		\begin{cases}
			\tfrac{k}{\alpha_{r}} \, \mathcal{F} [ U(\mathbf{r}^{\dagger}) ] 
			e^{\, k \, (z''_{0} - z^{\dagger}_{0})} & , 
			\quad 0 < k \le k_{\varepsilon} \: , \quad \left| \alpha_{r} \right| \ne 0 \\
			0 & , \quad k > k_{\varepsilon}
		\end{cases} \: .
		\label{eq:p-Fourier-min-phase}
	\end{equation}
	Note that $\mathcal{F} [ p(\mathbf{r}'') ]$ (equation \ref{eq:p-Fourier-min-phase}) does not depend on $\alpha_{d}$.
	As a consequence, changes in the differential operator $\alpha_{d}$, which carries all the information related to
	the observed potential field, do not affect $p(\mathbf{r}'')$ (equation \ref{eq:eqlayer-integral}).
	This concludes the proof.
\end{proof}

The representation of equivalent layers having minimum phase changes (Definition \ref{def:eql-min-phase}) in the
space domain is usually non-trivial and depends on the particular form of the real function $\alpha_{r}$ 
(equations \ref{eq:kernel-Fourier-min-phase} and \ref{eq:p-Fourier-min-phase}).
Such layers give raise to potential-field transformations (equation \ref{eq:transformed-field-integral-Fourier})
defined by $\alpha_{t} = \alpha_{r} \, \alpha_{f}$ (equation \ref{eq:alpha-t-generic}), which may also
have non-trivial representations in the space domain.

%As a direct consequence of Corollary \ref{corollary:minimum-phase-change-layers}, we have that any 
%equivalent layer having minimum phase changes (Definition \ref{def:eql-min-phase}) 
%simultaneously approximates the scalar potential $U$ 
%and its partial derivatives at any horizontal plane $z^{\dagger}_{0} < z^{t}_{0}$ above the 
%sources according to the error bound \ref{eq:epsilon-bound-data}.
%Consequently, any potential-field transformation 
%(equations \ref{eq:transformed-field-continuous-layer} and \ref{eq:transformed-field-integral-Fourier})
%does not depend on the particular observed potential field used to define the equivalent layer as well.
%This can be shown by noting that, for a particular $\alpha_{a} = \tilde{\alpha}_{r} \, \tilde{\alpha}_{d}$,
%we have $\tilde{\alpha}_{t} = \left(\nicefrac{\tilde{\alpha}_{r}}{\tilde{\alpha}_{\Phi}}\right) \, \alpha_{f}$ 
%(equation \ref{eq:alpha-t-generic}), which does not
%depend on the $\tilde{\alpha}_{d}$ associated with the particular observed potential field.

\subsection{Neumann's equivalent layers}

\begin{definition}
	An equivalent layer is called Neumann's equivalent layer if the Fourier transform of its physical-property distribution has 
	the form
	\begin{equation}
		\mathcal{F} [ p(\mathbf{r}'') ] =
		\begin{cases}
			k \, \mathcal{F} [ U(\mathbf{r}^{\dagger}) ] 
			e^{\, k \, (z''_{0} - z^{\dagger}_{0})} & , 
			\quad 0 < k \le k_{\varepsilon} \: , \\
			0 & , \quad k > k_{\varepsilon} \: ,
		\end{cases} 
		\label{eq:p-Fourier-Neumann}
	\end{equation}
	where $z''_{0} > z^{\dagger}_{0}$.
	\label{def:eql-Neumann}
\end{definition}
\begin{lemma}
	All Neumann's equivalent layers (Definition \ref{def:eql-Neumann}) have 
	minimum phase changes according to Definition \ref{def:eql-min-phase}.
	\label{lemma:eql-Neumann}
\end{lemma}
\begin{proof}
	From equation \ref{eq:p-Fourier}, we have that the solution $\mathcal{F} [ p(\mathbf{r}'') ]$ 
	given by equation \ref{eq:p-Fourier-Neumann} to define a Neumann's equivalent layer
	(Definition \ref{def:eql-Neumann}) is satisfied 
	only for the case in which $\alpha_{a} = \alpha_{d}$. This condition implies that 
	$\mathcal{F} [ p(\mathbf{r}'') ]$ is a particular case of equation \ref{eq:p-Fourier-min-phase},
	with $\alpha_{r} = 1$, and consequently has minimum phase changes
	according to Definition \ref{def:eql-min-phase}.
\end{proof}

\subsubsection{General aspects}

Let a given observed potential field $d^{o}(\mathbf{r})$ be reproduced by a 
Neumann's equivalent layer (Definition \ref{def:eql-Neumann}). 
In this case, the planar equivalent-layer integral (equation \ref{eq:eqlayer-integral}) assumes the form
\begin{equation}
	d^{o}(\mathbf{r}) \approx
	d(\mathbf{r})
	= \int\limits_{-\infty}^{\infty}\int\limits_{-\infty}^{\infty}
	p(\mathbf{r}'') \: D_{d} \Psi(\mathbf{r}, \mathbf{r}'') \:\: d\mathcal{S}''
	\: , \quad z < z''_{0} \: ,
	\label{eq:eqlayer-integral-Neumann-space}
\end{equation}
where $D_{d}$ is the same space-domain differential operator associated with
the observed potential field data $d^{o}(\mathbf{r})$ (equation \ref{eq:observed-data-space}).
As a consequence, Neumann's equivalent layers (Definition \ref{def:eql-Neumann}) preserve the physical
relationship between the observed potential field and the true sources.
It means that, if the observed potential field $d^{o}(\mathbf{r})$ is the $x$-component of the 
gravitational attraction due to the true sources, the predicted potential field $d(\mathbf{r})$
represents the $x$-component of the gravitational attraction due to the planar equivalent layer, for example.
The same is valid for any other potential field.

Another interesting aspect of Neumann's layers is that, given a particular physical-property distribution 
$\tilde{p}(\mathbf{r}'')$ satisfying \ref{eq:eqlayer-integral-Neumann-space}, potential-field transformations 
are defined in the Fourier domain (equation \ref{eq:transformed-field-integral-Fourier}) by using 
$\alpha_{t} = \alpha_{f}$ (equation \ref{eq:alpha-t-generic}). Hence, potential-field transformations
(equation \ref{eq:transformed-field-continuous-layer}) assume the following form in the space domain:
\begin{equation}
	f(\mathbf{r})
	= \int\limits_{-\infty}^{\infty}\int\limits_{-\infty}^{\infty}
	\tilde{p}(\mathbf{r}'') \: D_{f} \Psi(\mathbf{r}, \mathbf{r}'') \:\: d\mathcal{S}''
	\: , \quad z < z''_{0} \: ,
	\label{eq:transformed-field-integral-Neumann-space}
\end{equation}
where $D_{f}$ is the space-domain representation of the differential operator $\alpha_{f}$ 
(equations \ref{eq:diff-operators-generic}, \ref{eq:kernel-transformation-Fourier} 
and \ref{eq:alpha-t-generic}) associated with the desired transformed field $f(\mathbf{r})$.

\subsubsection{Direct relationship with the true sources}

It follows immediately from equation \ref{eq:p-Fourier-Neumann} that, in the space domain, the physical-property 
distribution $p(\mathbf{r}'')$ of Neumann's equivalent layers (equation \ref{eq:eqlayer-integral-Neumann-space})
approximates the vertical derivative
$\tfrac{\partial}{\partial z} U(\mathbf{r}'')$ (equation \ref{eq:potential-U-1st-deriv-space}) on the 
layer plane $z''_{0}$.
For the cases in which $d^{o}(\mathbf{r})$ represents gravity data, the physical-property 
distribution $p(\mathbf{r}'')$ approximates the gravity disturbance (or vertical component of the gravitational
attraction) due to the true sources on the layer plane.
If $d^{o}(\mathbf{r})$ represents magnetic data produced by sources with constant total magnetization direction, 
then $p(\mathbf{r}'')$ approximates the pseudogravity 
anomaly \citep[][ p. 345]{blakely1996} due to the true sources on the layer plane.
In this case, $p(\mathbf{r}'')$ is strictly positive because it is defined in terms of a strictly positive 
function $\sigma(\mathbf{r}')$ (equation \ref{eq:potential-U-space}) representing the 
total-magnetization intensity at all points $\mathbf{r}' = (x', y', z')^{\top} \in \mathcal{V}$ 
within the sources.
This result is in perfect agreement with that deduced by \citet{reis_etal2020} for
the particular case in which the observed potential field $d^{o}(\mathbf{r})$ is
the approximated total-field anomaly (equation \ref{eq:potential-U-oblique-2nd-deriv-space}).

Since $p(\mathbf{r}'')$ (equation \ref{eq:eqlayer-integral-Neumann-space}) approximates 
$\tfrac{\partial}{\partial z} U(\mathbf{r}'')$ (equation \ref{eq:potential-U-1st-deriv-space}) on the 
layer plane $z''_{0}$, then it shares some common properties with the true sources.
The first is the total mass, for gravity sources, or total dipole moment, for magnetic sources with
constant total-magnetization direction.
For the cases in which the layer plane is located above the sources, $p(\mathbf{r}'')$ (equation \ref{eq:eqlayer-integral-Neumann-space})
reproduces $\tfrac{\partial}{\partial z} U(\mathbf{r}'')$ (equation \ref{eq:potential-U-1st-deriv-space}) exactly 
(equations \ref{eq:lim-kepsilon-p-Fourier}--\ref{eq:lim-kepsilon-predicted-data-space}) 
and its integration on the layer plane gives the total mass (or total dipole moment) according to the
Gauss' Theorem \citep[e.g.,][]{hammer1945, lafehr1965, bhattacharyya1967}. 
Furthermore, \citet{medeiros-silva1995a} show that the horizontal first-order moments, second-order moment $xy$, 
difference between the second-order moments $xx$ and $yy$, as well as the horizontal coordinates of the sources center of mass
(or center of dipole moment) of the true sources can also be estimated from the $\tfrac{\partial}{\partial z} U(\mathbf{r}'')$
that they produce on a horizontal plane above them.
These quantities can only be approximated for the case in which the layer crosses the sources because, in this case,
$p(\mathbf{r}'')$ does not reproduce $\tfrac{\partial}{\partial z} U(\mathbf{r}'')$ exactly.
The important aspect to be emphasized here is that the horizontal distribution of the physical-property distribution $p(\mathbf{r}'')$
(equation \ref{eq:eqlayer-integral-Neumann-space}) on Neumann's equivalent layers is equal to or approximates 
that associated with the true sources.

\subsubsection{Implications for gravity-field transformations}

In the context of gravity exploration, Lemma \ref{lemma:eql-Neumann} gives the 
theoretical foundation for the common approach to jointly process gravity-gradient data 
\cite[e.g.,][]{barnes_lumley2011} or to convert gravity disturbance data into the 
gravity-gradient tensor \cite[e.g.,][]{piauilino_etal2019_seg, oliveirajr_etal2023}, for example.
Such gravity-data transformations usually presume, without any proof, the existence of a single layer of 
monopoles that satisfactorily reproduces all components of the gravitational field.
This assumption, however, is a direct consequence of 
Corollary \ref{corollary:minimum-phase-change-layers} and Lemma \ref{lemma:eql-Neumann}.

For example, let us consider the common two-step approach to convert gravity-disturbance into gravity
gradient data. The first step consists in solving the planar equivalent-layer integral (equation \ref{eq:eqlayer-integral})
numerically for obtaining an estimated physical-property distribution $\tilde{p}(\mathbf{r}'')$ by considering that the
observed potential field is given by $d^{o}(\mathbf{r}) = \tfrac{\partial}{\partial z} U(\mathbf{r})$ (equation \ref{eq:potential-U-1st-deriv-space})
and that the kernel is $a(\mathbf{r}, \mathbf{r}'') = \tfrac{\partial}{\partial z} \Psi(\mathbf{r}, \mathbf{r}'')$ 
(equation \ref{eq:Psi}).
The second step consists in using the estimated $\tilde{p}(\mathbf{r}'')$ to compute the gradient components.
The $uv$-component of the gradient tensor is computed by setting $D_{t} \equiv \tfrac{\partial^{2}}{\partial u \, \partial v}$
(equations \ref{eq:potential-U-2nd-deriv-space} and \ref{eq:transformed-field-integral-Neumann-space}),
where $u, v \in \{ x, y, z \}$.

Note that the first step consists in defining a Neumann's equivalent layer (Definition \ref{def:eql-Neumann}) with 
$\alpha_{d} = \alpha_{a} \equiv \alpha(1, \hat{\mathbf{z}})$ 
(equations \ref{eq:observed-data-Fourier}, \ref{eq:diff-operators-generic} and \ref{eq:kernel-Fourier}) 
which leads to a particular physical-property distribution 
$\tilde{p}(\mathbf{r}'')$ having a Fourier transform given by \ref{eq:p-Fourier-Neumann}. The second step is equivalent to 
compute the potential-field transformations (equation \ref{eq:transformed-field-integral-Fourier}) defined in the Fourier domain 
by using $\alpha_{t} = \alpha_{f}$ (equation \ref{eq:alpha-t-generic}), with 
$\alpha_{f} \equiv \alpha(1, \hat{\mathbf{u}}) \, \alpha(1, \hat{\mathbf{v}})$ (equation \ref{eq:diff-operators-generic}), 
where the unit vectors $\hat{\mathbf{u}}$ and $\hat{\mathbf{v}}$ are given by the canonical base vectors (equation \ref{eq:canonical-base-vector}).
This Fourier-domain operator $\alpha_{f}$ is equivalent to the space-domain operator $D_{t} \equiv D(1, \hat{\mathbf{u}}) \, D(1, \hat{\mathbf{v}})$
(equation \ref{eq:diff-operators-generic}).

\subsubsection{Implications for magnetic-field transformations}

\begin{lemma}
	A planar equivalent layer of dipoles defined in terms of the magnetic scalar potential or a 
	component of the magnetic induction 
	field produced along a given constant direction by sources with a constant total-magnetization direction
	is a Neumann's equivalent layer (Definition \ref{def:eql-Neumann}) 
	if and only if its constant magnetization direction is equal to that defining the total magnetization of the sources.
	\label{lemma:eql-Neumann-dipoles}
\end{lemma}
\begin{proof}
	Again, we get from equations \ref{eq:p-Fourier} and \ref{eq:p-Fourier-Neumann} that all Neumann's 
	equivalent layers must be defined by $\alpha_{a} = \alpha_{d}$.
	We split this proof into two parts.
	
	In the first part, we consider that the observed potential field is the magnetic scalar potential. 
	In this case, the differential operator 
	$\alpha_{d} \equiv \alpha(1, \hat{\mathbf{v}})$ (equation \ref{eq:diff-operators-generic}),
	where $\hat{\mathbf{v}}$ is a unit vector with the same direction as that defining the
	sources total magnetization. Consider that 
	$\alpha_{a} \equiv \alpha(1, \hat{\mathbf{h}})$ (equation \ref{eq:diff-operators-generic}),
	where $\hat{\mathbf{h}}$ is a unit vector different from $\hat{\mathbf{v}}$. Suppose for sake of
	contradiction that $\alpha(1, \hat{\mathbf{v}}) = \alpha(1, \hat{\mathbf{h}})$. Then, we use 
	equation \ref{eq:diff-operators-generic} and the fact that the dot product is distributive to obtain 
	$(\hat{\mathbf{v}} - \hat{\mathbf{h}})^{\top} \mathbf{k} = 0$, which means that $\mathtt{i} k_{x}$, 
	$\mathtt{i} k_{y}$ and $k = \sqrt{k_{x}^{2} + k_{y}^{2}}$ are linearly dependent, a contradiction.
	This implies that $\hat{\mathbf{h}} = \hat{\mathbf{v}}$, completing the first part of our proof.
	
	In the second part of our proof, we consider that the observed potential field is a component
	of the magnetic induction field along a given constant direction defined by the unit vector $\hat{\mathbf{u}}$.
	Now, we have that 
	$\alpha_{d} \equiv \alpha(1, \hat{\mathbf{u}}) \, \alpha(1, \hat{\mathbf{v}})$ and
	$\alpha_{a} \equiv \alpha(1, \hat{\mathbf{u}}) \, \alpha(1, \hat{\mathbf{h}})$ 
	(equation \ref{eq:diff-operators-generic}). We consider that these two differential operators
	are equal to each other for sake of contradiction, which also leads to
	$(\hat{\mathbf{v}} - \hat{\mathbf{h}})^{\top} \mathbf{k} = 0$, another contradiction.
	This implies again that $\hat{\mathbf{h}} = \hat{\mathbf{v}}$, completing the second part of our proof.
\end{proof}

In the context of magnetic exploration methods, Lemmas \ref{lemma:eql-Neumann} and 
\ref{lemma:eql-Neumann-dipoles} contribute to a deeper understanding about the use 
of dipole layers to perform magnetic-field transformations.
Surprisingly, most of the previously published studies about this topic are limited to
the use of a planar layer of dipoles as a purely mathematical strategy for processing total-field 
anomaly data, as we can see in the comments of \citet{pedersen1989} on the work of \citet{bhattacharyya_chan1977},
for example.
Few studies \cite[][]{pedersen1991, li_etal2014, reis_etal2020} have drawn
attention to the problem of proving the existence of a planar layer of dipoles
that reproduces total-field anomaly data.
%By Lemmas \ref{lemma:eql-Neumann} and \ref{lemma:eql-Neumann-dipoles},
%we obtain that planar layers of dipoles having a constant 
%magnetization direction equal to that associated with the total-magnetization of the true sources
%are (i) Neumann's equivalent layers (Definition \ref{def:eql-Neumann}), (ii) intrinsically related 
%to the true magnetic sources and (iii) directly connected to the solution of the Neumann's problem on 
%a plane (equation \ref{eq:U-Neumann}). It means that such layers do not have a purely mathematical,
%but also a physical meaning.

By combining Lemmas \ref{lemma:eql-Neumann} and \ref{lemma:eql-Neumann-dipoles} with
Corollary \ref{corollary:minimum-phase-change-layers}, we obtain that
Neumann's equivalent layers of dipoles having a constant magnetization direction 
provides the theoretical basis for the common two-step strategy to compute the RTP in the space domain, 
from total-field anomaly data, according to the pioneer work of \citet{silva1986}.

The first step consists in solving the planar equivalent-layer integral (equation \ref{eq:eqlayer-integral})
numerically by considering that (i) the observed potential field 
is satisfactorily described by the directional derivative $d^{o}(\mathbf{r}) = \hat{\mathbf{u}}^{\top} \mathbf{H}_{U}(\mathbf{r}) \, \hat{\mathbf{v}}$ 
(equation \ref{eq:potential-U-oblique-2nd-deriv-space}), where $\hat{\mathbf{u}}$ defines the constant direction of the main geomagnetic
field at the study area and $\hat{\mathbf{v}}$ the constant total-magnetization direction of the sources, and (ii)
the kernel is $a(\mathbf{r}, \mathbf{r}'') = \hat{\mathbf{u}}^{\top} \mathbf{H}_{\Psi}(\mathbf{r}, \mathbf{r}'') \, \hat{\mathbf{v}}$,
where $\mathbf{H}_{\Psi}(\mathbf{r}, \mathbf{r}'')$ is the Hessian matrix \citep[e.g.,][ p. 158]{kaplan2002} of
function $\Psi(\mathbf{r}, \mathbf{r}'')$ (equation \ref{eq:Psi}).
The second step consists in using the estimated $\tilde{p}(\mathbf{r}'')$ to compute the RTP anomaly 
$\tfrac{\partial^{2}}{\partial z^{2}} U(\mathbf{r})$ (see the subsection \textit{Some relevant potential fields for geophysical applications}).
This is made by evaluating the potential-field transformation (equation \ref{eq:transformed-field-continuous-layer}) 
with $t(\mathbf{r}, \mathbf{r}'') = \tfrac{\partial^{2}}{\partial z^{2}} \Psi(\mathbf{r}, \mathbf{r}'')$ 
(equation \ref{eq:Psi}).

Note that the first step consists in defining a Neumann's equivalent layer (Definition \ref{def:eql-Neumann}) with 
$\alpha_{d} = \alpha_{a} \equiv \alpha(1, \hat{\mathbf{u}}) \, \alpha(1, \hat{\mathbf{v}})$ (equations \ref{eq:observed-data-Fourier}, 
\ref{eq:diff-operators-generic} and \ref{eq:kernel-Fourier}), which leads to a dipole layer having the same constant 
magnetization direction of the true sources and a particular physical-property distribution 
$\tilde{p}(\mathbf{r}'')$ with Fourier transform given by \ref{eq:p-Fourier-Neumann}. The second step is equivalent to  
compute the potential-field transformation (equation \ref{eq:transformed-field-integral-Fourier}) defined in the Fourier domain 
by using $\alpha_{t} = \alpha_{f}$ (equation \ref{eq:alpha-t-generic}), with 
$\alpha_{f} \equiv \alpha(2, \hat{\mathbf{z}})$ (equations \ref{eq:diff-operators-generic} and \ref{eq:canonical-base-vector}).
In this case, it can be verified from equations \ref{eq:diff-operators-generic} and \ref{eq:transformed-field-integral-Fourier3} that the 
transformed potential field $f(\mathbf{r}^{\ddagger})$ approximates the RTP anomaly $\tfrac{\partial^{2}}{\partial z^{2}} U(\mathbf{r}^{\ddagger})$ 
on the plane $z^{\ddagger}_{0}$.

It is important to underline two significant aspects that stands out from our theoretical approach to compute the RTP 
in terms of a Neumann's equivalent layer.
The first is that it has never been clearly defined in the pioneer work of \citet{silva1986}, nor in the succeeding studies 
\citep[e.g., ][]{leao-silva_1989, guspi_novara2009, oliveirajr_etal2013, li_etal2014, reis_etal2020, takahashi-etal2022} 
using the equivalent-layer technique. 
We do not presume, but mathematically prove the existence of a dipole layer that allows computing the RTP.
It is important to stress that \citet{silva1986} cites the work of \citet{roy1962} to justify that ``the magnetic 
field due to arbitrary sources can be exactly fitted by the magnetic field due to a horizontal equivalent
layer of dipoles''. 
Based on this, \citet{silva1986} proposed a method that (i) estimates the physical-property distribution on a dipole layer
having a constant magnetization direction equal to that of the true sources and (ii) sets the magnetization of this 
layer in the vertical direction to compute the RTP.
This two-step approach, that is widely accepted by the scientific community and we know to be correct from practice, is an empirical 
generalization of the \citeauthor{roy1962}'s (\citeyear{roy1962}) work.
\citeauthor{roy1962} did not discuss the existence of a dipole layer that fits the observed magnetic-field data,
but presented a general outline to deduce the solutions of the Dirichlet's and Neumann's problems on a plane 
(equation \ref{eq:U-Dirichlet-Neumann}).

The second significant aspect of our approach is that it also offers 
the theoretical link between the RTP in the space domain, via equivalent-layer technique, and 
in the Fourier domain, via classical approach \citep[e.g., ][ p. 331]{blakely1996}. 
This can be seen by first noting that the common two-step approach to compute the RTP via equivalent-layer technique is 
equivalent to a single-step approach defined by equation
\ref{eq:transformed-field-integral-Fourier-single-step2}, with $\alpha_{f} = \alpha(2, \hat{\mathbf{z}})$ 
(equations \ref{eq:diff-operators-generic}, \ref{eq:canonical-base-vector}, 
\ref{eq:transformed-field-integral-Fourier} and \ref{eq:alpha-t-generic}) and 
$\alpha_{d} \equiv \alpha(1, \hat{\mathbf{u}}) \, \alpha(1, \hat{\mathbf{v}})$ (equations \ref{eq:observed-data-Fourier} and 
\ref{eq:diff-operators-generic}), 
where $\hat{\mathbf{u}}$ defines the constant direction of the main geomagnetic
field at the study area and $\hat{\mathbf{v}}$ the constant total-magnetization direction of the sources.
Then, by considering that the plane $z^{\dagger}_{0}$ containing the observed data $d^{o}(\mathbf{r}^{\dagger})$ coincides with the 
plane $z^{\ddagger}_{0}$ where the transformed field $f(\mathbf{r}^{\ddagger})$ is computed, equation \ref{eq:transformed-field-integral-Fourier-single-step2} 
reduces to that commonly used in the literature \citep[e.g., ][ p. 331]{blakely1996} to perform the RTP in the 
Fourier domain.
This result gives the reason by which the equivalent-layer technique ultimately suffers the same instability as 
the Fourier-domain RTP operator at low latitudes, as have already been pointed out by \citet{li_etal2014}.


\subsection{Dipole layers with arbitrary constant magnetization direction}

Consider now that (i) the observed potential field 
is satisfactorily described by the directional derivative $d^{o}(\mathbf{r}) = \hat{\mathbf{u}}^{\top} \mathbf{H}_{U}(\mathbf{r}) \, \hat{\mathbf{v}}$ 
(equation \ref{eq:potential-U-oblique-2nd-deriv-space}), where $\hat{\mathbf{u}}$ defines the constant direction of the main geomagnetic
field at the study area and $\hat{\mathbf{v}}$ the constant total-magnetization direction of the sources, and (ii)
the kernel is $a(\mathbf{r}, \mathbf{r}'') = \hat{\mathbf{u}}^{\top} \mathbf{H}_{\Psi}(\mathbf{r}, \mathbf{r}'') \, \hat{\mathbf{h}}$,
where $\mathbf{H}_{\Psi}(\mathbf{r}, \mathbf{r}'')$ is the Hessian matrix \citep[e.g.,][ p. 158]{kaplan2002} of
function $\Psi(\mathbf{r}, \mathbf{r}'')$ (equation \ref{eq:Psi}) and $\hat{\mathbf{h}}$ is a constant unit vector not necessarily equal to 
$\hat{\mathbf{v}}$.
In this case, we get from equation \ref{eq:p-Fourier} that a solution of the planar equivalent layer 
can be obtained by using 
$\alpha_{d} \equiv \alpha(1, \hat{\mathbf{u}}) \, \alpha(1, \hat{\mathbf{v}})$ and 
$\alpha_{a} \equiv \alpha(1, \hat{\mathbf{u}}) \, \alpha(1, \hat{\mathbf{h}})$ 
(equation \ref{eq:diff-operators-generic}), so that the physical-property distribution assumes the following form 
in the Fourier domain:
\begin{equation}
	\mathcal{F} [ p(\mathbf{r}'') ] =
	\begin{cases}
		k \, \tfrac{\left(\mathbf{k}^{\top} \hat{\mathbf{v}} \right)}{\left(\mathbf{k}^{\top} \hat{\mathbf{h}} \right)} \,
		\mathcal{F} \left[ U(\mathbf{r}^{\dagger}) \right] 
		e^{\, k \, (z''_{0} - z^{\dagger}_{0})} & , 
		\quad 0 < k \le k_{\varepsilon} \\
		0 & , \quad k > k_{\varepsilon} 
	\end{cases} \: .
	\label{eq:p-Fourier-dipoles}
\end{equation}
This particular solution (equation \ref{eq:p-Fourier-dipoles}) was deduced by considering
that $\hat{\mathbf{u}}$ is a unit vector having the same constant direction of the main geomagnetic field at the study area.
However, it can be easily verified that this solution can also be deduced by considering that $\hat{\mathbf{u}}$ is one of the
canonical base vectors $\hat{\mathbf{x}}$, $\hat{\mathbf{y}}$ or $\hat{\mathbf{z}}$ (equation \ref{eq:canonical-base-vector}).
In this case, the observed potential field represents a Cartesian component of the magnetic induction field.
Additionally, we also deduce the same solution given by equation \ref{eq:p-Fourier-dipoles} if we consider that the 
observed potential field is the magnetic scalar potential. In this case, $\alpha_{d}$ and $\alpha_{a}$ are defined
without the term $\alpha(1, \hat{\mathbf{u}})$.
Hence, the planar equivalent layer formed by dipoles with physical-property distribution 
given by equation \ref{eq:p-Fourier-dipoles} simultaneously explain the magnetic scalar potential and
magnetic field components.

Another interesting aspect to be emphasized on the particular solution defined by equation \ref{eq:p-Fourier-dipoles} is that,
according to Lemma \ref{lemma:eql-Neumann-dipoles}, it 
does not necessarily lead to equivalent layers having minimum phase changes (Definition \ref{def:eql-min-phase}) owing to the complex ratio 
$\nicefrac{(\mathbf{k}^{\top} \hat{\mathbf{v}})}{(\mathbf{k}^{\top} \hat{\mathbf{h}})}$.
For the particular case in which $\hat{\mathbf{h}}$ coincides with the constant magnetization direction $\hat{\mathbf{v}}$ of the true sources,
then equation \ref{eq:p-Fourier-dipoles} reduces to that defining the physical-property distribution of a Neumann's layer 
(equation \ref{eq:p-Fourier-Neumann}), which has minimum phase changes according to Definition \ref{def:eql-min-phase}.

Note that the physical-property distribution $\mathcal{F} [ p(\mathbf{r}'') ]$ (equation \ref{eq:p-Fourier-dipoles}) associated 
with an arbitrarily magnetized dipole layer leads to potential-field transformations (equations 
\ref{eq:transformed-field-continuous-layer} and \ref{eq:transformed-field-integral-Fourier}) with 
$\alpha_{t}$ (equation \ref{eq:alpha-t-generic}) defined in terms of the complex ratio 
$\nicefrac{\alpha_{a}}{\alpha_{d}}$, where 
$\alpha_{d} \equiv \alpha(1, \hat{\mathbf{u}}) \, \alpha(1, \hat{\mathbf{v}})$ and 
$\alpha_{a} \equiv \alpha(1, \hat{\mathbf{u}}) \, \alpha(1, \hat{\mathbf{h}})$ 
(equation \ref{eq:diff-operators-generic}).
It means, for example, that the physical-property distribution $\mathcal{F} [ p(\mathbf{r}'') ]$ (equation \ref{eq:p-Fourier-dipoles}) 
cannot be used to compute the RTP by using the common approach with $\alpha_{t} = \alpha(2, \hat{\mathbf{z}})$ 
(equations \ref{eq:diff-operators-generic}, \ref{eq:canonical-base-vector} and \ref{eq:alpha-t-generic}).
This common approach can only be used if the constant magnetization direction $\hat{\mathbf{h}}$ coincides with 
that of the true sources $\hat{\mathbf{v}}$.
This fact is widely known by the geophysical community in practice and has long been assumed to be true
without any mathematical proof by several authors 
\citep[e.g., ][]{silva1986, leao-silva_1989, guspi_novara2009, oliveirajr_etal2013, takahashi-etal2022}. 
Our results offer the first theoretical explanation to this fact.

\subsection{Alternative equivalent layers}

\subsubsection{``Generalized equivalent sources'' proposed by \citet{cordell1992}}

Let the planar equivalent-layer integral (equation \ref{eq:eqlayer-integral}) be numerically solved by
using a kernel function $a(\mathbf{r}, \mathbf{r}'') = \Psi(\mathbf{r}, \mathbf{r}'')$ (equation \ref{eq:Psi}). 
Note that, in this case, the predicted potential field $d(\mathbf{r})$ 
represents the gravitational potential due to the equivalent layer.
The solution $p(\mathbf{r})$ obtained by this approach is defined from equation \ref{eq:p-Fourier}, in the Fourier domain, 
by setting $\alpha_{a} = 1$ (equation \ref{eq:diff-operators-generic}):
\begin{equation}
	\mathcal{F} [ p(\mathbf{r}'') ] =
	\begin{cases}
		k \, \alpha_{d} \, \mathcal{F} \left[ U(\mathbf{r}^{\dagger}) \right] 
		e^{\, k \, (z''_{0} - z^{\dagger}_{0})} & , 
		\quad 0 < k \le k_{\varepsilon} \\
		0 & , \quad k > k_{\varepsilon} 
	\end{cases} \: .
	\label{eq:p-Fourier-generalized-eqs}
\end{equation}
It is also important noting that 
this particular solution (equation \ref{eq:p-Fourier-generalized-eqs}) approximates the vertical derivative of the observed 
potential field on the layer plane and results in potential-field transformations (equations \ref{eq:transformed-field-continuous-layer} and 
\ref{eq:transformed-field-integral-Fourier}) defined by $\alpha_{t} = \nicefrac{\alpha_{f}}{\alpha_{d}}$ (equation \ref{eq:alpha-t-generic}).

This strategy, called \textit{generalized equivalent sources}, was proposed by \citet{cordell1992} and recently adapted by 
\citet{soler_uieda2021} to perform gravity-data processing. 
Those authors considered only the particular case in which observed potential field is gravity disturbance data 
$d^{o}(\mathbf{r}) = \tfrac{\partial}{\partial z} U(\mathbf{r})$ (equation \ref{eq:observed-data-space}), 
with $\alpha_{d} \equiv \alpha(1, \hat{\mathbf{z}})$ (equations \ref{eq:observed-data-Fourier}, 
\ref{eq:diff-operators-generic} and \ref{eq:canonical-base-vector}).
In this particular case, the solution given by equation \ref{eq:p-Fourier-generalized-eqs} 
(i) approximates $\tfrac{\partial^{2}}{\partial z^{2}} U(\mathbf{r})$ (equation \ref{eq:potential-U-2nd-deriv-space}) 
in the space domain, 
(ii) gives rise to equivalent layers having minimum phase changes (Definition \ref{def:eql-min-phase}), with 
$\alpha_{r} = \nicefrac{1}{k}$ (equation \ref{eq:p-Fourier-min-phase}), and 
(iii) results in potential-field transformations defined by $\alpha_{t} = \nicefrac{\alpha_{f}}{k}$ (equation \ref{eq:alpha-t-generic}).

\subsubsection{Alternative approach proposed by \citet{guspi_novara2009}}

\citet{guspi_novara2009} proposed an alternative two-step approach to compute the RTP via equivalent-layer technique. 
The first step consists in solving the planar 
equivalent-layer integral (equations \ref{eq:eqlayer-integral} and \ref{eq:eqlayer-integral-Fourier}) numerically 
by considering that (i) the observed potential field 
is satisfactorily described by the directional derivative 
$d^{o}(\mathbf{r}) = \hat{\mathbf{u}}^{\top} \mathbf{H}_{U}(\mathbf{r}) \, \hat{\mathbf{v}}$ 
(equation \ref{eq:potential-U-oblique-2nd-deriv-space}), where $\hat{\mathbf{u}}$ defines the constant direction of the main geomagnetic
field at the study area and $\hat{\mathbf{v}}$ the constant total-magnetization direction of the sources, and (ii)
the kernel is given by 
\begin{equation}
	a(\mathbf{r}, \mathbf{r}'') = \tfrac{\left( z''_{0} - z \right)}{2 \, \pi} \, 
	\ln{\Big[} \left( z''_{0} - z \right) + \| \mathbf{r}'' - \mathbf{r} \| {\Big]} - 
	\tfrac{\| \mathbf{r}'' - \mathbf{r} \|}{2 \, \pi} \: ,
	\label{eq:kernel-Fourier-Guspi}
\end{equation}
which represents the double integral of the function $\Psi(\mathbf{r}, \mathbf{r}'')$ (equation \ref{eq:Psi}) 
along the vertical direction $z$.
The second step of the alternative approach proposed by \citet{guspi_novara2009} computes the RTP anomaly 
by setting the potential-field transformation (equation \ref{eq:transformed-field-continuous-layer}) with 
$t(\mathbf{r}, \mathbf{r}'') = \Psi(\mathbf{r}, \mathbf{r}'')$ (equation \ref{eq:Psi}).

Note that the first step consists in obtaining an equivalent layer defined by 
$\alpha_{d} = (\mathbf{k}^{\top} \hat{\mathbf{u}}) (\mathbf{k}^{\top} \hat{\mathbf{v}})$ (equations \ref{eq:observed-data-Fourier} 
and \ref{eq:diff-operators-generic}) and 
$\alpha_{a} = \nicefrac{(\mathbf{k}^{\top} \hat{\mathbf{u}}) (\mathbf{k}^{\top} \hat{\mathbf{h}})}{k^{2}}$, 
where the term $\nicefrac{1}{k^{2}}$ represents the double integration along the vertical direction. In this case,  
the physical-property distribution $\mathcal{F} [ p(\mathbf{r}'') ]$ (equation \ref{eq:p-Fourier}) assumes the form:
\begin{equation}
	\mathcal{F} [ p(\mathbf{r}'') ] =
	\begin{cases}
		k^{3} \, \tfrac{\left(\mathbf{k}^{\top} \hat{\mathbf{v}} \right)}{\left(\mathbf{k}^{\top} \hat{\mathbf{h}} \right)} \,
		\mathcal{F} \left[ U(\mathbf{r}^{\dagger}) \right] 
		e^{\, k \, (z''_{0} - z^{\dagger}_{0})} & , 
		\quad 0 < k \le k_{\varepsilon} \\
		0 & , \quad k > k_{\varepsilon} 
	\end{cases} \: .
	\label{eq:p-Fourier-guspi}
\end{equation}
The second step is equivalent to compute the potential-field transformation (equation \ref{eq:transformed-field-integral-Fourier}) 
defined in the Fourier domain by using $\alpha_{t} = 1$ (equation \ref{eq:alpha-t-generic}).
In this case, it can be verified that the transformed potential field $\mathcal{F} [ f(\mathbf{r}^{\ddagger}) ]$ 
(equation \ref{eq:transformed-field-integral-Fourier}) actually approximates the RTP anomaly for the particular case 
in which $\mathcal{F} [ p(\mathbf{r}'') ]$ (equation \ref{eq:p-Fourier-guspi}) is defined by using $\hat{\mathbf{h}} = \hat{\mathbf{v}}$, 
i.e., for the case in which the constant magnetization direction of the equivalent layer coincides with that of the true sources.
This result shows the effectiveness of our approach in providing the theoretical background not only for classical 
equivalent-layer methods, but also for unconventional ones.